\begin{abstract}
LLM Agents now route user tasks through tool servers, model APIs and other Agents over encrypted channels, and a fast-growing literature shows that the resulting traffic reveals queries, task types and Agent identities. Defenses are starting to appear, but they are evaluated with methods inherited from website fingerprinting, which assume that cover traffic is free of side effects and that one link is the unit of protection. We systematize fifteen Agent studies (forty-four settings) against twenty-two classical traffic-analysis anchors and find that both assumptions fail in Agent execution. First, tool feedback re-enters model requests, so the same task is exposed on several independently controlled links: a defense that protects one link must be judged against every view its attacker retains. Second, Agent cover is often \emph{executed} as application work, so a privacy result is incomplete without the original task's outcome under the same intervention. The literature has not yet closed this gap. None of the four mapped query/workflow privacy settings reports both an adapted attacker and matched task utility, and the published research-decoy result pairs its privacy gain with utility measured under a different intervention. A paired MCP experiment with two LLM planners shows why the omission matters. Four cover semantics leave recognition within two of 120 tasks for one planner, yet differ by up to 108 of 120 fixtures in business-state integrity. Under padded metadata, residual recognition follows the part of the trace that the cover schedule leaves unprotected, not how deeply the cover executes. We contribute an access-aware attack and defense map, a classical-to-Agent transfer matrix, a four-step evaluation procedure and a source-dependency analysis showing which conclusions survive without preprints.
\end{abstract}
